\section{Related Work}
\label{sec:related}

Recent work has shown that different fairness notions are often
mutually
exclusive~\cite{KleinbergMR16,Chouldechova17,FriedlerSV16}. Unsurprisingly
then, different fairness notions have corresponded to different
algorithms and optimization frameworks.  Previously introduced
fairness notions have generally split along several axis:
classification vs. regression and individual vs. group fairness and
disparate treatment.  Most of previous work has focused on
classification, despite the ubiquity of regression in real world
applications with fairness concerns.

In classification, one line of work aims to achieve the group fairness
notion known as \emph{statistical parity}, i.e. to avoid disparate
impact (see e.g. \cite{PedreshiRT08, KamiranCP10, CaldersV10,
  KamiranC11, LuongRT11, KamiranKZ12, KamishimaAAS12, HajianD13,
  FeldmanFMSV15,FishKL16,AdlerFFRSSV16}).  Statistical parity requires
a predictor to predict each label at similar rates across different
groups.  This definition can be at odds with accuracy especially when
the two groups are inherently different.  \citet{HardtPS16} introduced
a new notion of group fairness called \emph{equality of odds},
partially to alleviate this friction, and partially arguing that
equality of odds more accurately captured what it would mean for a
classifier to be equally ``good'' for two groups. Equality of odds has
a very intuitive interpretation for classification: it requires
similarity of misclassification rates across groups (rather than
forcing the marginal classification rates in the two groups to
coincide). Optimizing for accuracy subject to an equality of odds
constraint was recently shown to be NP-hard \cite{WoodworthGOS17};
work following this result presented efficient heuristics for the
problem \cite{ZafarVGG17, WoodworthGOS17}. We also study
fairness definitions which we can implement efficiently but our interest is in studying
the trade-off between fairness and accuracy and the relationship between different notions
of fairness.

Although equality of odds is also defined for regression, it is very
difficult to determine empirically whether a regressor's output is
conditionally independent of the protected attribute (conditioned on
the true label), as each true label may be seen only once.


\citet{CaldersKKAZ13} introduced the study of statistical parity's
analog in regression settings, (called \emph{equal means} and
\emph{balanced residuals}). More recently \citet{JohnsonFS16} also
studied fairness for regression problems and formalized several
notions for \emph{impartial estimates} based on the causal
relationship between \emph{sensitive attributes}, \emph{legitimate
  attributes}, \emph{suspect attributes} and the label.  Both groups
consider group fairness; our group fairness notion differs from these
as we incorporate the similarity of pairs (through the function $d$)
in our definition though the specific choice of
$d(y,y')=c$ for some $c\in\R$ and all $y$ and $y'$ would recover equal means.

To achieve any of these fairness notions, one needs to decide whether
or not to allow for \emph{disparate treatment} (allowing for different
treatment, or different models, for different groups)\footnote{Any
  classifier that uses sensitive attributes in its decision making is
  implicitly fitting separate models to the two populations, and while
  this might seem unfair, it has been argued that it is actually
  necessary for fairness (most notably in \citet{DworkHPRZ12}).}  , and
where in the learning process to enforce fairness: preprocessing of
data (e.g.~\cite{KamiranC11}); inprocessing, during the training of a
model as either a constraint or incorporated into the objective
function (e.g.~\cite{FishKL16, KamishimaAAS12,ZafarVGG15}); or
postprocessing, where data is labeled by some black-box model and then
relabeled as a function only of the original labels
(e.g. ~\cite{HardtPS16}).  Our approach in this paper falls into the
in-processing category, by encoding fairness as a regularizer (an
approach previously studied in
e.g.~\cite{KamishimaAAS12,ZafarVGG15,ZafarVGG17}). Our work differs
from previous work in several aspects by primarily focusing on
regression, and that our family of fairness measures draws inspiration
from the idea that \emph{similar} instances should be treated
\emph{similarly}~\cite{DworkHPRZ12, ZemelWSPD13}.

% Finally, recent papers have proven several impossibility results for
% satisfying different fairness notions
% simultaneously~\cite{KleinbergMR16,FriedlerSV16}, necessitating that
% the algorithm designer should commit to a particular definition of
% fairness.


% \sj{moved from intro}
%
% \textit{Individual fairness} definitions are criteria which bind on
% the classifier between pairs of individuals, whereas \textit{group
% fairness} is a constraint on the aggregate performance of the model
% (e.g. equalizing the marginal distribution of positive examples in
% two different populations).  We take inspiration from
% \cite{DworkHPRZ12} that fairness should mean that ``similar people
% are treated similarly". \cite{DworkHPRZ12} assumed that a similarity
% metric between individuals was known; here, we follow
% \citet{JosephKMR16} and \citet{HardtPS16} and say that the metric is
% simply defined based on the true labels.
%
%But even with this heuristic in hand, previous works on fairness take very different approaches to the following natural questions:
%\begin{itemize}
%\item Should the requirement that ``similar people be treated similarly" bind at the individual or group level?
%\item Can disparate treatment be naturally incorporated into the model, and what are the consequences of doing so?
%\item Do fairness definitions that are tailored to classification naturally extend to the regression setting?
%\item How does the fairness vs. accuracy tradeoff vary across datasets? Is there practical guidance about the source of variation?
%\end{itemize}

\iffalse
Ultimately as a data analyst, whether a fairness notion is needed that emphasizes group or individual fairness, fits separate models
to separate groups or uses one model, and applies to classification or regression, depends on  (often qualitative)
attributes of the specific problem at hand, and the goals and responsibilities of the algorithm designer. In light of this fact, a framework for
fairness that allows the algorithm designer to make these choices in a unified and consistent way is desirable.
In this paper we posit such a framework  that applies to regression or classification, has natural variants that interpolate between
notions of group and individual fairness, easily incorporates disparate treatment if desired, and is efficiently implementable.
We then exploit the tractability of our framework to present a comparative
study of the tradeoffs between fairness and accuracy across a variety
of fairness notions, models, and datasets. We draw our datasets from
the high stakes domains of criminal justice, lending, admissions. Interestingly,
the tradeoff between fairness and accuracy varies widely across datasets, while
the effects of \textit{disparate treatment}, or fitting one vs. two models are relatively muted.
This indicates that toy examples and moral arguments aside, with real datasets the decision of
whether to model different  groups separately may not be an important one.
\fi
%\sn{Put in the above intro section / sales pitch. As I was writing it
%I started to feel like the main contribution of this paper is its flexibility as a means
%of trading off different fairness goals and assumptions.
%Feel free to tear it apart / make the results more specific, make the references more accurate and inclusive }

\iffalse This has the potential to both improve the overall utility of
the fair classifier, while avoiding undesirable properties such as
rejecting more qualified black applicants than qualified white
applicants, for example. ~\citet{HardtPS16} showed how to achieve
these notions via post-processing of a (possibly unfair) black-box
classifier, and prove that the resulting classifier is optimal under
certain assumptions. Even more recently \citet{WoodworthGOS17} show
that a post-processing approach, as suggested by \citet{HardtPS16},
can be arbitrarily suboptimal in derived settings.  They present a
nearly-optimal statistical procedure for achieving equality of odds,
and argue that the associated computational problem is
intractable. They then suggest a second moment relaxation of the
non-discrimination definition, for which learning is tractable for
linear predictors.

They also note that this predictor can be derived from the optimal
least squares predictor in the sense of \citet{HardtPS16}.  The paper
of \citet{WoodworthGOS17} opens several interesting practical
questions concerning learning equality of odds predictors on real
datasets:
\begin{itemize}
\item What other tractable relaxations of learning classifiers
  satisfying equality of odds are there, and how do they perform
  empirically?
\item While post-processing can be arbitrarily bad on contrived examples, how does it compare in practice to learning with a fairness
constraint?
\end{itemize}

The above questions serve as motivation for the definitions and results below.

\citet{ZafarVGG17}
\sj{to do: discuss how we came up with these notions of fairness and
  how they would compare with previously studied notions of fairness.}
 \fi 